\documentclass[11pt]{article}
\usepackage{mathunal-base}
\mathunalfuente{serif}

\renewcommand{\examtitulo}{Primer Examen Parcial de Álgebra Lineal}
\renewcommand{\examfecha}{9 de Marzo del 2015}

\begin{document}

\examheader

\vspace{4pt}
\noindent
\begin{minipage}[t]{0.68\linewidth}
\textbf{Puntaje. Sólo para uso Oficial}\\[4pt]
\begin{tabular}{|c|c|c|c|c|c|c|}
\hline
1--6 & 7 & 8 & 9 & 10 & \textbf{TOTAL} & \textbf{NOTA} \\ \hline
 & & & & & & \\ \hline
\end{tabular}
\end{minipage}

\vspace{6pt}
\noindent\textbf{Instrucciones:} La duración del examen es de 1 hora y 50 minutos. El examen consta de diez preguntas en dos hojas impresas por ambos lados, verifique que su examen esté completo y consérvelo con el gancho. En las preguntas con procedimiento justifique sus respuestas en los espacios asignados. No está permitido sacar hojas en blanco ni ningún tipo de apuntes durante el examen, verifique que su celular esté apagado. No se permite el uso de calculadora.

\vspace{6pt}
\noindent\textbf{IDENTIFICACIÓN}

\vspace{8pt}
\noindent Nombre: \dotfill\ Cédula \dotfill

\vspace{4pt}
\noindent Profesor: \dotfill\ Grupo \dotfill

\vspace{10pt}
\begin{center}\textbf{I. Completación}\end{center}
\noindent En las preguntas 1 a 6 complete los espacios en blanco. \textbf{NOTA}: En esta sección se califica sólo la respuesta y no se tiene en cuenta el procedimiento.

\begin{enumerate}
\item{}[8pt] Determine todos los escalares $k$ tales que $||k\mathbf{v}||=3$ si $\mathbf{v}=[-1,2,7,-3]$, $k=$ \dotfill

\item{}[8pt] Determine si el sistema homogéneo tiene soluciones no triviales (Escriba SI o NO en la línea debajo de cada sistema)

\vspace{6pt}
\noindent
\begin{minipage}[t]{0.5\linewidth}
$\begin{aligned}
x_1+3x_2+5x_3+x_4 &=0\\
4x_1-7x_2-3x_3-x_4 &=0\\
3x_1+2x_2+7x_3+8x_4 &=0
\end{aligned}$

\vspace{4pt}
\dotfill
\end{minipage}%
\begin{minipage}[t]{0.46\linewidth}
$\begin{aligned}
x+2y+3z &=0\\
y+4z &=0\\
5z &=0
\end{aligned}$

\vspace{4pt}
\dotfill
\end{minipage}

\item{}[6pt] Sean $A,B$ matrices $4\times 5$, $C,D$ y $E$ de tamaños $5\times 2$, $4\times 2$ y $5\times 4$ respectivamente. Determine si las siguientes matrices están o no bien definidas (escriba SI o NO, según el caso).

\vspace{6pt}
\noindent (a) $BA$ \dotfill

\vspace{4pt}
\noindent (b) $AC+D$ \dotfill

\vspace{4pt}
\noindent (c) $AE+B$ \dotfill

\item{}[6pt] Si una matriz no nula de tamaño $3\times 3$, tiene todas sus entradas iguales entonces su rango es \dotfill .

\item{}[6pt] Si $A$ y $B$ son matrices simétricas, escriba SI o NO, según el caso.

\vspace{6pt}
\noindent (a) $A^T$ \dotfill\ es simétrica.

\vspace{4pt}
\noindent (b) $BA$ \dotfill\ es simétrica.

\vspace{4pt}
\noindent (c) $AA^T$ \dotfill\ es simétrica.

\item{}[8pt] Sea $T = \left\{\begin{bmatrix}1\\1\\1\end{bmatrix},\begin{bmatrix}0\\1\\1\end{bmatrix},\begin{bmatrix}1\\0\\k\end{bmatrix}\right\}$

\vspace{4pt}
\noindent Determine los valores de $k$ para que $T$ sea L.I. $k=$ \dotfill
\end{enumerate}

\vspace{6pt}
\begin{center}\textbf{II. Solución con Procedimiento}\end{center}

\begin{enumerate}
\setcounter{enumi}{6}
\item{}[10pt] Determine si los vectores generan a $\mathbb{R}^3$
\[
\begin{bmatrix}1\\1\\1\end{bmatrix},\ \begin{bmatrix}2\\2\\0\end{bmatrix},\ \begin{bmatrix}3\\0\\0\end{bmatrix},\ \begin{bmatrix}1\\-1\\1\end{bmatrix}
\]

\vspace{8cm}

\item{}[26pt] Un biólogo ha colocado tres cepas bacterianas (denotadas como I, II y III) en un tubo de ensayo, donde serán alimentadas con tres distintas fuentes alimenticias (A, B y C). Cada día 400 unidades de A, 500 de B y 300 de C se colocan en el tubo de ensayo, y cada bacteria consume cierto número de unidades de cada alimento por día, como se muestra en la siguiente tabla. ¿Cuántas bacterias de cada cepa pueden coexistir en el tubo de ensayo y consumir todo el alimento?.

\vspace{6pt}
\begin{center}
\begin{tabular}{|l|c|c|c|}
\hline
Alimento & Cepa bacteriana I & Cepa bacteriana II & Cepa bacteriana III \\ \hline
Alimento A & 1 & 2 & 0 \\ \hline
Alimento B & 2 & 1 & 3 \\ \hline
Alimento C & 1 & 1 & 1 \\ \hline
\end{tabular}
\end{center}

\begin{enumerate}
\item{} Plantee un sistema de ecuaciones lineales que permita resolver el problema. Defina claramente las variables a utilizar.
\item{} Encuentre un intervalo, para cada variable libre, donde las soluciones tienen sentido.
\item{} Si se tiene la cantidad máxima de bacterias de la cepa III, ¿Cuántas bacterias de cada cepa pueden haber?
\end{enumerate}

\vspace{6cm}

\item{}[12pt] Sea $H = gen\left(\begin{bmatrix}1\\-2\\1\\2\end{bmatrix},\begin{bmatrix}0\\1\\-1\\-1\end{bmatrix}\right)$. Dé un sistema de ecuaciones lineales cuyo conjunto de soluciones sea $H$. En otras palabras, determine a $H$ por medio de restricciones en las entradas.

\vspace{8cm}

\item{}[10pt] Sean $v$ y $w$ vectores en $\mathbb{R}^n$. Demostrar que $\{v, v+w, v-w\}$ es un conjunto linealmente dependiente.

\vspace{4cm}
\end{enumerate}

\examfooter
\end{document}
